% -----------------------------------------------------------------------
% REWRITTEN related work — added VQA foundations, query difficulty,
% multimodal RAG, and FrugalGPT cascade distinction; all existing
% paragraphs expanded
% -----------------------------------------------------------------------

\section{Extended Related Work}\label{sec:related_ext}

\noindent\textbf{Visual Question Answering.}
VQA~\cite{antol2015vqa} introduced the task of answering free-form natural language questions about images.
VQA-v2~\cite{goyal2017vqav2} reduced language-prior biases via complementary image pairs; GQA~\cite{hudson2019gqa} introduced compositional reasoning over scene graphs; TextVQA~\cite{singh2019textvqa} highlighted OCR-dependent visual understanding.
Analysis of VQA performance has consistently found that question type is strongly predictive of model accuracy and failure modes: counting and text-recognition questions are universally harder, while yes/no questions are systematically easier~\cite{kafle2017analysis, agrawal2018dont}.

\noindent\textbf{Pre-Retrieval Query Difficulty Estimation.}
Estimating query difficulty before retrieval---from query features alone---is a longstanding problem in information retrieval~\cite{zhao2008predicting, hauff2008improved, mothe2005linguistic}.
These methods use lexical, syntactic, and statistical properties of the query to predict retrieval quality without executing the retrieval step, enabling cost-aware allocation of retrieval resources.
VOILA adapts this paradigm to multimodal perception: we estimate, from question-linguistic features alone, the probability that a VLM succeeds at each visual fidelity---without executing the VLM or accessing any image.
This makes VOILA a direct generalization of pre-retrieval query difficulty estimation to the cross-modal, cost-sensitive fidelity selection setting.

\noindent\textbf{FrugalGPT-Style Cascades and Confidence-Based Routing.}
FrugalGPT~\cite{chen2023frugalgpt} reduces LLM inference cost by routing queries through a cascade of models in ascending cost order, using the cheapest model that produces a sufficiently confident answer~\cite{shnitzer2023large, wang2024routellm}.
Adapted to fidelity selection, this yields a natural \emph{sequential cascade baseline}: run the VLM at the lowest fidelity; if output log-probability falls below a threshold, escalate to the next fidelity; repeat until confident.
VOILA differs from this paradigm in two fundamental ways.
First, the cascade is \emph{post-retrieval}---it incurs retrieval and VLM inference costs at every intermediate fidelity before reaching the final answer, meaning it cannot avoid the costs it aims to minimize.
Second, cascades rely on the VLM's own confidence as the escalation signal, which Observation 2 shows is unreliable when visual context is insufficient: the VLM produces confident-but-wrong answers at low fidelity, so the cascade does not escalate when it should.
VOILA avoids both problems by selecting the target fidelity \emph{before any retrieval} using question-conditioned prediction, eliminating intermediate costs and bypassing post-hoc confidence signals entirely.
We compare directly against a log-prob cascade baseline in \S\ref{sec:results}.

\noindent\textbf{Adaptive Inference and Model Routing.}
Beyond FrugalGPT, a broad line of work studies adaptive computation for NLP: early-exit transformers~\cite{zhou2020bert, elbayad2020depth} skip layers for ``easy'' inputs; CALM~\cite{schuster2022confident} exits generation early when token confidence is sufficient; RouteLLM~\cite{wang2024routellm} and hybrid routing~\cite{ding2024hybrid} dispatch queries across model tiers based on predicted difficulty.
All of these methods \emph{assume the input context is fully available} and optimize the computation applied to it.
VOILA is orthogonal: it optimizes \emph{what input context to acquire} before any computation begins, and can be composed with model routing to jointly optimize retrieval and inference costs.

\noindent\textbf{Adaptive Visual Resolution and Concurrent Work.}
A concurrent line of work addresses adaptive resolution selection for visual inputs.
CARES~\cite{kimhi2025cares} selects the minimal sufficient input resolution for a given image-query pair before forwarding to the target VLM, using a compact VLM to extract query-conditioned features from the image.
Prompt-guided prefiltering~\cite{azizian2026promptguided} identifies spatially task-relevant regions within an already-retrieved image before compression.
Both approaches reduce VLM compute or transmission cost, but both still require \emph{image access} to make the resolution or region decision.
VOILA is distinguished by its strictly pre-retrieval setting: in tiered-storage RAG systems, the image has not been fetched when the routing decision must be made.
VOILA therefore operates exclusively from question text---a strictly harder setting that eliminates retrieval cost at the source rather than reducing processing cost afterward.
Post-retrieval token compression~\cite{weng2024learning} and lightweight cross-modal fusion~\cite{farazi2021accuracy} reduce the cost of \emph{processing} a retrieved image; \textsc{VOILA} is complementary, deciding which representation to acquire in the first place.
% TODO before submission: verify kimhi2025cares (arxiv:2510.19496) and azizian2026promptguided (arxiv:2604.00314) are in example_paper.bib

\noindent\textbf{Multimodal Retrieval-Augmented Generation.}
Retrieval-augmented generation (RAG)~\cite{lewis2020rag} extends language models with retrieved context; multimodal extensions retrieve images, video frames, or memory traces~\cite{yasunaga2023retrievalaugmented, chen2022murag}.
In agentic memory settings~\cite{park2023generative, wang2024locomo}, agents must decide which stored memories to retrieve given a query.
All existing multimodal RAG frameworks treat retrieval as binary (retrieve or not) and assume a fixed representation quality for whatever is retrieved.
\textsc{VOILA} introduces a finer-grained retrieval decision: \emph{at what fidelity to retrieve}.

\noindent\textbf{Value-of-Information and Cost-Sensitive Acquisition.}
The value-of-information (VOI) framework~\cite{bilgic2011value} formalizes the expected utility of acquiring additional evidence given current beliefs.
Budgeted feature acquisition~\cite{karkkainen2019cost} and active acquisition in multimodal settings~\cite{kossen2022active} apply VOI to tabular and sequential data.
VOILA is the first to apply VOI to pre-retrieval fidelity selection in visual question answering, where the ``features'' are not tabular attributes but image representations at different compression levels, and the action space is structured by a physical retrieval cost function over storage tiers.

% -----------------------------------------------------------------------
% OLD RELATED WORK (kept for reference)
% -----------------------------------------------------------------------
\begin{comment}
\noindent\textbf{Adaptive Inference and Model Routing (old).}
Prior work studies adaptive computation by routing queries across models or dynamically allocating compute, including early-exit and adaptive depth mechanisms \cite{zhou2020bert, schuster2022confident} and cost-aware LLM routing \cite{chen2023frugalgpt, shnitzer2023large, ding2024hybrid, wang2024routellm}.
These methods optimize \emph{how} inference is performed given fixed inputs, whereas VOILA focuses on \emph{what information} to retrieve before inference.

\noindent\textbf{Systems Context for Multimodal Retrieval (old).}
Multimodal systems commonly store visual data across heterogeneous tiers, motivating retrieval-aware optimization.
Related work examines tiered storage and retrieval costs in cloud and mobile systems \cite{hort2021survey, zhang2024autonomous}, as well as agentic memory and long-horizon multimodal reasoning \cite{park2023generative, wang2024locomo}.
VOILA explicitly exploits this structure by selecting the minimum sufficient fidelity before accessing higher-cost tiers.
\end{comment}
